% GENERATED FILE. Edit canonical lessons, then run npm run generate:books.
% canonical-source-hash: fnv1a64:090ca9277d6d2cac
% canonical-lessons: FR-C354-cv, FR-C354-candidat, FR-C354-candidature, FR-C354-postuler, FR-C354-embauche

\chapter{Je cherche du travail}
\label{ch:fr-je-cherche-du-travail}
\hlchaptermodality{354}

\begin{chapteropening}
\textbf{By the end of this chapter:} \emph{I can apply for a job: send my CV and my application, say who the candidates are and talk about being taken on.}

Complete the last lesson of chapter 354: I can apply for a job: send my CV and my application, say who the candidates are and talk about being taken on.
\end{chapteropening}

\section[le CV]{le CV — the CV, the résumé}

\label{lesson:FR-C354-cv}



\begin{quote}
Before the new one: say the French for to lessen, then the French for to make up for.
\end{quote}

\subsection*{You'll want to know: le CV}
\textbf{le CV} — \textquotedblleft{}the CV, the résumé\textquotedblright{}.

\begin{quote}
J'ai envoyé mon CV à dix entreprises.
\end{quote}

\begin{quote}
\emph{I sent my CV to ten companies.}
\end{quote}

Said \emph{cé-vé}: short for Latin \emph{curriculum vitae}. Masculine.

\begin{cousinweb}
\emph{CV} is short for \emph{curriculum vitae}, Latin for \textquotedblleft{}a running of life\textquotedblright{}: \emph{curriculum} \textquotedblleft{}a running, a course\textquotedblright{}, from \emph{currere} \textquotedblleft{}to run\textquotedblright{}, and \emph{vīta} \textquotedblleft{}life\textquotedblright{} (Lewis \& Short).
\end{cousinweb}

\begin{grammarlens}[title={What you've built}]
One more word for telling, explaining and describing at B1.
\end{grammarlens}

\subsection*{Your turn}
\begin{itemize}
\raggedright
  \item \emph{Say it:} \emph{le CV}
  \item \emph{Say it:} \emph{le CV}, once more
  \item \emph{Say it:} a sentence of your own with \emph{le CV}
  \item \emph{Recall:} say the French for to lessen, then the French for to make up for, then say \emph{le CV} again
\end{itemize}

\begin{tcolorbox}[breakable,colback=teal!4,colframe=teal!35!black,title={Before you move on}]
What does le CV mean? (\textquotedblleft{}The CV, the résumé\textquotedblright{}.) Say it once more.
\end{tcolorbox}
\section[le candidat]{le candidat — the candidate, the applicant}

\label{lesson:FR-C354-candidat}



\begin{quote}
Before the new one: say the French for to make up for, then the French for the CV.
\end{quote}

\subsection*{You'll want to know: le candidat}
\textbf{le candidat} — \textquotedblleft{}the candidate, the applicant\textquotedblright{}.

\begin{quote}
Il y a vingt candidats pour ce poste.
\end{quote}

\begin{quote}
\emph{There are twenty applicants for this job.}
\end{quote}

Feminine \emph{candidate}. Also in elections.

\begin{cousinweb}
\emph{candidat} comes from Latin \emph{candidātus} \textquotedblleft{}clothed in white\textquotedblright{} (EtymDB; Lewis \& Short), from \emph{candidus} \textquotedblleft{}dazzling white\textquotedblright{}: in Rome, people seeking office wore a whitened toga.
\end{cousinweb}

\begin{grammarlens}[title={What you've built}]
One more word for telling, explaining and describing at B1.
\end{grammarlens}

\subsection*{Your turn}
\begin{itemize}
\raggedright
  \item \emph{Say it:} \emph{le candidat}
  \item \emph{Say it:} \emph{le candidat}, once more
  \item \emph{Say it:} a sentence of your own with \emph{le candidat}
  \item \emph{Recall:} say the French for to make up for, then the French for the CV, then say \emph{le candidat} again
\end{itemize}

\begin{tcolorbox}[breakable,colback=teal!4,colframe=teal!35!black,title={Before you move on}]
What does le candidat mean? (\textquotedblleft{}The candidate, the applicant\textquotedblright{}.) Say it once more.
\end{tcolorbox}
\section[la candidature]{la candidature — the application (for a job)}

\label{lesson:FR-C354-candidature}



\begin{quote}
Before the new one: say the French for the CV, then the French for the candidate.
\end{quote}

\subsection*{You'll want to know: la candidature}
\textbf{la candidature} — \textquotedblleft{}the application (for a job)\textquotedblright{}.

\begin{quote}
J'ai envoyé ma candidature hier soir.
\end{quote}

\begin{quote}
\emph{I sent my application last night.}
\end{quote}

\emph{Poser sa candidature} — to apply.

\begin{cousinweb}
\emph{candidature} is built from \emph{candidat} and the suffix \emph{-ure}. \emph{candidat} comes from Latin \emph{candidātus} \textquotedblleft{}clothed in white\textquotedblright{}, from \emph{candidus} \textquotedblleft{}white\textquotedblright{} (Lewis \& Short).
\end{cousinweb}

\begin{grammarlens}[title={What you've built}]
One more word for telling, explaining and describing at B1.
\end{grammarlens}

\subsection*{Your turn}
\begin{itemize}
\raggedright
  \item \emph{Say it:} \emph{la candidature}
  \item \emph{Say it:} \emph{la candidature}, once more
  \item \emph{Say it:} a sentence of your own with \emph{la candidature}
  \item \emph{Recall:} say the French for the CV, then the French for the candidate, then say \emph{la candidature} again
\end{itemize}

\begin{tcolorbox}[breakable,colback=teal!4,colframe=teal!35!black,title={Before you move on}]
What does la candidature mean? (\textquotedblleft{}The application (for a job)\textquotedblright{}.) Say it once more.
\end{tcolorbox}
\section[postuler]{postuler — to apply (for a job)}

\label{lesson:FR-C354-postuler}



\begin{quote}
Before the new one: say the French for the candidate, then the French for the application.
\end{quote}

\subsection*{You'll want to know: postuler}
\textbf{postuler} — \textquotedblleft{}to apply (for a job)\textquotedblright{}.

\begin{quote}
Elle a postulé à un poste de vendeuse.
\end{quote}

\begin{quote}
\emph{She applied for a job as a shop assistant.}
\end{quote}

\emph{Postuler à} or \emph{pour un poste}.

\begin{cousinweb}
\emph{postuler} is a learned borrowing from Latin \emph{postulāre} \textquotedblleft{}to ask, to demand\textquotedblright{} (EtymDB; Lewis \& Short).
\end{cousinweb}

\begin{grammarlens}[title={What you've built}]
One more word for telling, explaining and describing at B1.
\end{grammarlens}

\subsection*{Your turn}
\begin{itemize}
\raggedright
  \item \emph{Say it:} \emph{postuler}
  \item \emph{Say it:} \emph{postuler}, once more
  \item \emph{Say it:} a sentence of your own with \emph{postuler}
  \item \emph{Recall:} say the French for the candidate, then the French for the application, then say \emph{postuler} again
\end{itemize}

\begin{tcolorbox}[breakable,colback=teal!4,colframe=teal!35!black,title={Before you move on}]
What does postuler mean? (\textquotedblleft{}To apply (for a job)\textquotedblright{}.) Say it once more.
\end{tcolorbox}
\section[l'embauche]{l'embauche — the hiring, being taken on}

\label{lesson:FR-C354-embauche}



\begin{quote}
Before the new one: say the French for the application, then the French for to apply.
\end{quote}

\subsection*{You'll want to know: l'embauche}
\textbf{l'embauche} — \textquotedblleft{}the hiring, being taken on\textquotedblright{}.

\begin{quote}
Après l'entretien, il a eu une promesse d'embauche.
\end{quote}

\begin{quote}
\emph{After the interview, he got a written job offer.}
\end{quote}

\emph{Une promesse d'embauche} — a written job offer. Feminine.

\begin{cousinweb}
\emph{embauche} is formed from the verb meaning to take on a worker (EtymDB), built on an old word \emph{bauche} \textquotedblleft{}workshop\textquotedblright{}. The origin of \emph{bauche} is not settled.
\end{cousinweb}

\begin{grammarlens}[title={What you've built}]
One more word for telling, explaining and describing at B1.
\end{grammarlens}

\subsection*{Your turn}
\begin{itemize}
\raggedright
  \item \emph{Say it:} \emph{l'embauche}
  \item \emph{Say it:} \emph{l'embauche}, once more
  \item \emph{Say it:} a sentence of your own with \emph{l'embauche}
  \item \emph{Recall:} say the French for the application, then the French for to apply, then say \emph{l'embauche} again
\end{itemize}

\begin{tcolorbox}[breakable,colback=teal!4,colframe=teal!35!black,title={Before you move on}]
What does l'embauche mean? (\textquotedblleft{}The hiring, being taken on\textquotedblright{}.) Say it once more.
\end{tcolorbox}
